\documentclass[10pt,a4paper]{amsart}
\usepackage[margin=19mm]{geometry}
\usepackage{amsmath,amssymb,mathtools}
\usepackage[scr=rsfs,cal=euler]{mathalfa}
\usepackage{xfrac}
\usepackage[unicode,hidelinks,bookmarks=false]{hyperref}
\newcommand{\ie}{i.e.\ }
\newcommand{\cf}{cf.\ }
\newcommand{\resp}{respectively\ }
\newcommand{\Subsection}{\subsection}
\providecommand{\nobreakdash}{\nobreak\hbox{-}}
\setlength{\emergencystretch}{2em}
\multlinegap0pt
\usepackage{bm}
\usepackage{bbm}
\usepackage{dsfont}
\usepackage{xspace}
\usepackage{cancel}
\usepackage{mathtools}
\usepackage{amsthm}
\makeatletter
\@ifundefined{thm}{\newtheorem{thm}{Theorem}}{}
\makeatother
\newcommand{\C}{{\mathbb C}}
\title[Adventures in algebraic geometry]{Adventures in algebraic geometry}
\author{Lev Borisov}
\date{Source: arXiv:2407.10041v2 (CC BY 4.0)}
\begin{document}
\maketitle

\noindent\emph{Source and licence.} Adventures in algebraic geometry. Version arXiv:2407.10041v2 (CC BY 4.0), distributed under the Creative Commons Attribution 4.0 International licence, as recorded in the arXiv licence metadata for this exact version and shown on the abstract page https://arxiv.org/abs/2407.10041v2. Full rights notice: licenses/arxiv-2407.10041-CC-BY-4.0.txt.\par\smallskip

\noindent\textit{Editorial scope.} Counted unit: the Thm whose source label is \texttt{XG} in the article (source0.tex lines 2596--2602, source label \texttt{XG}), delivered together with the article's own complete original proof of it (source0.tex lines 2604--2665, one \texttt{proof} environment of 62 lines). No mathematics is written, repaired or re-derived: statement and proof are the article's own text line for line. Required context retained uncounted: none. Every definition, display and label of those lines is retained unchanged. Omitted from this package and not redistributed: 1--2595, 2603--2603, 2666--5533. Nothing inside a retained excerpt is omitted or altered: every retained body is delivered exactly as the article's lines are written, so no omission receipt of any kind accompanies this package. Citations: the retained excerpts carry no citation command at all, so no bibliography entry of the article is transcribed and no reference is redistributed. Reference adaptation: none was needed and none was made; every reference inside the delivered unit resolves inside it, so no unresolved reference or citation is produced. Printed slips kept literal: no printed word, symbol, hypothesis or number of the counted statement or of its proof was altered. Layout and class: the article's own class is \texttt{article} with packages including amsmath, amsthm, amsfonts, amssymb, amscd, lastpage, enumerate, fancyhdr, mathrsfs, xcolor, graphicx, listings, hyperref, enumitem; the wrapper is the lane's pinned \texttt{amsart} editorial wrapper at 10pt on A4, which restates the theorem-like environments the retained text needs and adds the article's own macro definitions that the retained text needs (\texttt{\textbackslash C}), with extra wrapper packages (bm, bbm, dsfont, xspace, cancel, mathtools). The delivered numbering is the unit's own. Assets: no retained span contains a figure environment, a graphics-inclusion command, a PDF-page inclusion, a table or a TikZ picture; no image file is copied into this package, the package declares no asset dependency and no image file is redistributed with it. Licence: version arXiv:2407.10041v2 of this article is distributed under the Creative Commons Attribution 4.0 International licence, as recorded in the arXiv licence metadata for this exact version and shown on the abstract page https://arxiv.org/abs/2407.10041v2. A full rights notice is delivered at licenses/arxiv-2407.10041-CC-BY-4.0.txt. Characters: every retained excerpt is pure ASCII, so the delivered engine, XeLaTeX with its default Latin Modern text font, reports no missing character.
\begin{thm}\label{XG}
The set $X/G$ has a natural structure of an affine algebraic variety. More precisely, the subring of $G$-invariant elements
$$
(\C[x_1,\ldots,x_n]/I)^G
$$
is a finitely generated $\C$-algebra, and its maximal ideals are in a natural bijection with $G$-orbits on $X$.
\end{thm}

\begin{proof}
Consider the elements $y_{ik} = g_i(x_k)$ for all $g_i\in G$ and all $1\leq k\leq n$. We can use this set of generators to write $X$ as an algebraic subvariety of $\C^{n|G|}$, so that the action of $G$ comes from a permutation action on the coordinates of $\C^n$.
Then we have 
$$
(\C[x_1,\ldots,x_n]/I)^G \cong (\C[y_1,\ldots,y_l]/J)^G\cong \C[y_1,\ldots,y_l]^G/J^G.
$$
In the second isomorphism we used the fact that every element $r$ of the invariant ring  $(\C[y_1,\ldots,y_l]/J)^G$
can be lifted to an element $\hat r$ of  $\C[y_1,\ldots,y_l]$ and then 
$$\frac 1{|G|}\sum_{g\in G} g(\hat r)=r\,{\rm mod}\,J.$$
Thus, to prove that $(\C[x_1,\ldots,x_n]/I)^G$ is a finitely generated $\C$-algebra, it suffices to prove it for
$ \C[y_1,\ldots,y_l]^G$, with the permutation action of $G$.
We actually don't care too much that the action is a permutation action, but we will use a weaker statement that the action of 
$G$ on $A= \C[y_1,\ldots,y_l]$ preserves the natural grading on $A$ by the total degree in $y_i$.  

\smallskip
Consider the ideal $A A_+^G\subset A$ which is generated by all homogeneous elements of $A^G$ of positive degree. Since $A$ is Noetherian, we can find a finite set of generators $f_1,\ldots, f_k$ of $A A_+^G$, and we can even assume that all $f_i$ are $G$-invariant homogeneous elements of $A$. We will now prove that
$f_i$ generate the ring $A^G$ as a $\C$-algebra, by induction on degree.

\smallskip
In degree zero, we have $A_0 = A^G_{0}=\C$, and there is nothing to prove. Suppose we have proved that every homogeneous element of $A^G$ of degree at most $d$ can be written as a polynomial in $f_1,\ldots, f_k$ with constant coefficients. Then for a homogeneous element $f\in A^G$ of degree $d+1$ we use $f\in AA_+^G$ to write 
$$
f = \sum_{i=1}^k a_i f_i
$$
where $a_i$ are some elements of $A$ of degree $d+1-\deg f_i$. We will now use the averaging trick
\begin{align*}
f = \frac 1{|G|} \sum_{g\in G}g(f) =\frac 1{|G|} \sum_{i=1}^k  \sum_{g\in G}g(a_i f_i) 
=\frac 1{|G|} \sum_{i=1}^k  \sum_{g\in G}g(a_i) f_i
\\
=\sum_{i=1}^k  f_i \left(\frac 1{|G|} \sum_{g\in G}g(a_i) \right) = \sum_{i=1}^k f_i h_i
\end{align*}
where $h_i$ are $G$-invariant elements of degree $d+1-\deg f_i$. By induction assumption, $h_i$ can be written as polynomials in $f_1,\ldots,f_k$, therefore so can be $f$.

\smallskip
Now that we have established that $(\C[x_1,\ldots,x_n]/I)^G$ is finitely generated, let us prove the bijection between its maximal ideals and the set of orbits of $G$-action on the set of maximal ideals of $\C[x_1,\ldots,x_n]/I$.\footnote{ Recall that points in $X$ correspond to maximal ideals in $R=\C[x_1,\ldots,x_n]/I$ by Hilbert's Nullestellensatz.} In one direction, to any maximal ideal $m\subset R$ we associate the ideal $m^G = m\cap R^G\subset R^G$. The ideal $m^G$ is maximal because for any 
$x\in R^G\setminus m^G$ there exists $y\in R$ such that $xy=1\,{\rm mod}\,m$. Then
the element 
$$
\prod_{g\in G} (1 - g(x y))  = \prod_{g\in G} (1 - x g(y))  = 1 - x\hat y
$$
(with $G$-invariant $\hat y$) lies in $m^G$, which shows that $x \,{\rm mod}\,m^G$ is invertible. It is clear that different $m$ in the same $G$-orbit give the same $m^G$.

\smallskip
In the other direction, for a maximum ideal $I\subset R^G$ consider the ideal $RI$ that it generates in $R$. If $RI$ contains $1$
then there holds
$$
1 = \sum_i r_i x_k
$$
with $x_k\in I$. Averaging gives $1 =\sum_i \left (\frac 1{|G|}r_i\right)x_k\in I,$ contradiction. Therefore, $RI$ is contained in some maximal ideal $m$. It is clear that $m^G$ contains $I$ and is therefore equal to it, since $I$ is maximal.

\smallskip
It now suffices to show that if $m_1$ and $m_2$ satisfy $m_1^G=m_2^G$, then $m_1$ and $m_2$ are in the same $G$-orbit.
Consider $I_1 = \bigcap_{g\in G} g(m_1)$ and $I_2= \bigcap_{g\in G} g(m_2)$. If $I_1+ I_2$ were contained in a maximum ideal $m$,
we would have 
$$
m\supseteq \bigcap_{g\in G} g(m_1),~m\supseteq \bigcap_{g\in G} g(m_2).
$$
If $m$ is not equal to any $g(m_1)$, then there are elements $x_g\in g(m_1),x_g\not\in m$, and their product is in $ \bigcap_{g\in G} g(m_1)$ but not in $m$. Similarly, $m$ must be equal to one of $g(m_2)$, and this is impossible, since $m_1$ and $m_2$ are in different $g$-orbits. Therefore, $I_1+I_2=R$ and\footnote{ You have been warned.}
$$
1 = a_1+a_2,~a_i\in I_i.
$$
Moreover, since $I_i$ are $G$-invariant, when we average over $G$, we may assume that $a_i$ are $G$-invariant. However, this means that $a_i\in m_i^G$, which contradicts $1\not\in m_1^G=m_2^G$.
\end{proof}

\end{document}
